% !TeX root = main.tex
\newpage
\section{VerAForceDynamics Source Flow}

\newcommand{\VerAForceFunctionSection}[1]{%
  \subsection{VerAForceDynamics:#1}%
  \label{vera-force:#1}%
}

\texttt{VerAForceDynamics} treats overlap area as an instantaneous central
repulsive force field.  One source-particle invocation detects all current
particle and wall contacts, immediately vector-sums their forces into a
source-local \texttt{totalForce}, and writes only the source velocity and
position.  The implementation is shaped to transfer directly to one GLSL
compute invocation per source particle.  This model is the preserved
baseline: it uses physical particle radii and physical wall-ghost geometry for
starting overlaps, with no effective starting-contact state.

\begin{figure}[ht]
\centering
\resizebox{\textwidth}{!}{\begin{tikzpicture}[
    node distance=0.27in and 0.44in,
    function/.style={
        draw=black!60,
        fill=black!3,
        rounded corners=2pt,
        minimum width=2.18in,
        minimum height=0.28in,
        align=center,
        font=\ttfamily\small
    },
    entry/.style={
        function,
        fill=green!8,
        draw=green!40!black,
        line width=0.65pt
    },
    wall/.style={
        function,
        fill=blue!7,
        draw=blue!45!black
    },
    reused/.style={
        draw=orange!55!black,
        fill=orange!10,
        rounded corners=2pt,
        minimum width=1.26in,
        minimum height=0.22in,
        inner xsep=2pt,
        align=center,
        font=\ttfamily\scriptsize
    },
    flow/.style={-{Stealth[length=5pt]}, draw=black!65, line width=0.45pt},
    reusecall/.style={-{Stealth[length=4pt]}, draw=orange!60!black, dashed, line width=0.4pt},
    loop/.style={flow, rounded corners=0pt}
]

\newcommand{\VerAForceFlowNode}[2]{\hyperref[vera-force:#1]{\strut #2}}

\node[entry] (main) {main(SourceID)};

% Particle contacts are processed immediately when detected.
\node[function, below left=0.46in and 0.72in of main] (particle-scan)
    {particle cell scan};
\node[function, below=of particle-scan] (particle-process)
    {ProcessParticleCollision};
\node[function, below=of particle-process] (init)
    {\VerAForceFlowNode{InitializeContactState}{InitializeContactState}};
\node[function, below=of init] (state)
    {\VerAForceFlowNode{GetContactState}{GetContactState}};
\node[function, below=of state] (append)
    {\VerAForceFlowNode{AppendContactSlot}{AppendContactSlot}};
\node[function, below=of append] (geometry)
    {\VerAForceFlowNode{GetParticleGeometry}{GetParticleGeometry}};
\node[function, below=of geometry] (particle-force)
    {AccumulateContactForce};

% Each configured wall is processed exactly once.
\node[wall, below right=0.46in and 0.72in of main] (wall-scan)
    {four-wall scan};
\node[wall, below=of wall-scan] (wall-process)
    {ProcessWallCollision};
\node[wall, below=of wall-process] (wall-init)
    {InitializeWallContactState};
\node[wall, below=of wall-init] (wall-geometry)
    {GetWallGhostGeometry};
\node[wall, below=of wall-geometry] (wall-append)
    {AppendWallContactSlot};
\node[wall, below=of wall-append] (wall-force)
    {AccumulateContactForce};

% Source-local totalForce is complete before integration.
\node[entry, below=0.58in of main |- particle-force.south] (velocity)
    {CalcVelocity(SourceID, totalForce)};
\node[entry, below=of velocity] (position)
    {CalcPosition(SourceID)};

% Reused helpers are repeated beside their callers.
\node[reused, right=0.18in of append] (append-slots)
    {\VerAForceFlowNode{GetContactSlots}{GetContactSlots}};
\node[reused, right=0.18in of geometry] (geometry-position)
    {\VerAForceFlowNode{GetParticlePosition}{GetParticlePosition}};
\node[reused, right=0.18in of particle-force] (particle-stiffness)
    {GetContactStiffness};
\node[reused, left=0.18in of wall-geometry] (wall-position)
    {\VerAForceFlowNode{GetParticlePosition}{GetParticlePosition}};
\node[reused, left=0.18in of wall-append] (wall-slots)
    {\VerAForceFlowNode{GetContactSlots}{GetContactSlots}};
\node[reused, left=0.18in of wall-force] (wall-stiffness)
    {GetContactStiffness};
\node[reused, left=of velocity] (velocity-mass)
    {\VerAForceFlowNode{GetParticleMass}{GetParticleMass}};
\node[reused, left=of velocity-mass] (velocity-start)
    {\VerAForceFlowNode{GetStartFrameVelocity}{GetStartFrameVelocity}};
\node[reused, left=of position] (position-current)
    {\VerAForceFlowNode{GetParticlePosition}{GetParticlePosition}};

\draw[flow] (main) -- (particle-scan);
\draw[flow] (main) -- (wall-scan);

\draw[flow] (particle-scan) -- node[right,font=\scriptsize]{contact found} (particle-process);
\draw[flow] (particle-process) -- (init);
\draw[flow] (init) -- (state);
\draw[flow] (state) -- (append);
\draw[flow] (append) -- (geometry);
\draw[flow] (geometry) -- (particle-force);
\draw[loop] (particle-force.west) -- ++(-0.20in,0)
    |- node[pos=0.12,left,font=\scriptsize]{next detected contact} (particle-scan.west);

\draw[flow] (wall-scan) -- node[right,font=\scriptsize]{wall 1--4} (wall-process);
\draw[flow] (wall-process) -- (wall-init);
\draw[flow] (wall-init) -- (wall-geometry);
\draw[flow] (wall-geometry) -- (wall-append);
\draw[flow] (wall-append) -- (wall-force);
\draw[loop] (wall-force.east) -- ++(0.20in,0)
    |- node[pos=0.12,right,font=\scriptsize]{next wall} (wall-scan.east);

\draw[flow] (particle-force.south) -- ++(0,-0.18in) -| (velocity.north);
\draw[flow] (wall-force.south) -- ++(0,-0.18in) -| (velocity.north);
\draw[flow] (velocity) -- (position);

\draw[reusecall] (append) -- (append-slots);
\draw[reusecall] (geometry) -- (geometry-position);
\draw[reusecall] (particle-force) -- (particle-stiffness);
\draw[reusecall] (wall-geometry) -- (wall-position);
\draw[reusecall] (wall-append) -- (wall-slots);
\draw[reusecall] (wall-force) -- (wall-stiffness);
\draw[reusecall] (velocity) -- (velocity-mass);
\draw[reusecall] (velocity-mass) -- (velocity-start);
\draw[reusecall] (position) -- (position-current);

\node[font=\sffamily\scriptsize, text=black!65, above=0.08in of particle-scan]
    {particle contacts};
\node[font=\sffamily\scriptsize, text=blue!55!black, above=0.08in of wall-scan]
    {wall contacts};

\end{tikzpicture}

}
\caption{Current VerAForceDynamics source-owned force flow}
\end{figure}

\clearpage

\VerAForceFunctionSection{CollisionRun}

\texttt{CollisionRun} is the Python harness for the shader-shaped execution
schedule:

\begin{verbatim}
snapshot frame-start positions and velocities
initialize totalForce for every source
for each SourceID:
    detect particle contacts
        ProcessParticleCollision(TargetID, SourceID, totalForce)
for each SourceID:
    for each configured wall:
        ProcessWallCollision(SourceID, wall, totalForce)
capture frame-start pair diagnostics
for each SourceID:
    CalcVelocity(SourceID, totalForce)
for each SourceID:
    CalcPosition(SourceID)
snapshot end positions for diagnostics
\end{verbatim}

The Python harness uses a direct particle scan where GLSL uses cell occupancy
lists.  Both schedules complete one source-local force sum before integration.
The end-position snapshot is diagnostic only and is not a second force scan.

\VerAForceFunctionSection{ProcessParticleCollision}

\texttt{ProcessParticleCollision(TargetID, SourceID, totalForce)} initializes
the current contact geometry, records a diagnostic contact slot, and
immediately adds the contact force to the source-local \texttt{totalForce}.
It does not write the target particle.

\VerAForceFunctionSection{ProcessWallCollision}

\texttt{ProcessWallCollision(SourceID, wall, totalForce)} creates the
configured wall ghost geometry and immediately adds its force to
\texttt{totalForce}.  Each of the four configured walls is scanned exactly
once for each source particle.

\VerAForceFunctionSection{InitializeWallContactState}

\texttt{InitializeWallContactState(SourceID, wall)} represents a wall contact
with a stationary ghost having the same radius and mass as the source
particle.  The contact normal is perpendicular to the wall.  The configurable
\texttt{wall\_contact\_offset} controls how far the ghost center lies beyond
the wall boundary.

\VerAForceFunctionSection{InitializeContactState}

\texttt{InitializeContactState(SourceID, TargetID)} obtains the current
particle-particle geometry and records one source-owned diagnostic contact
slot.

\VerAForceFunctionSection{GetContactState}

\texttt{GetContactState(SourceID, TargetID)} obtains a reusable current-frame
contact slot.  The slot contains geometry and diagnostic values only; it is
not used later to calculate force.

\VerAForceFunctionSection{AppendContactSlot}

\texttt{AppendContactSlot(SourceID, TargetID)} records an active
particle-contact slot owned by the source.

\VerAForceFunctionSection{GetContactSlots}

\texttt{GetContactSlots(SourceID)} returns the fixed diagnostic contact-slot
array owned by the source particle.

\VerAForceFunctionSection{GetParticleGeometry}

\texttt{GetParticleGeometry(SourceID, TargetID)} calculates overlap area,
center distance, penetration depth, and the source-to-target normal from the
immutable frame-start position snapshot.

\VerAForceFunctionSection{GetParticlePosition}

\texttt{GetParticlePosition(ParticleID)} returns a particle position from the
selected immutable position snapshot.

\VerAForceFunctionSection{GetStartFrameVelocity}

\texttt{GetStartFrameVelocity(ParticleID)} returns the immutable frame-start
velocity used by integration and relative-normal-velocity diagnostics.

\VerAForceFunctionSection{AccumulateContactForce}

For every current particle or wall contact,
\texttt{AccumulateContactForce(SourceID, contact\_state, totalForce)} applies:

\begin{equation}
F_{ij}=q_{ij}A_{ij},
\end{equation}

\begin{equation}
\mathbf F_i \mathrel{+}= -F_{ij}\mathbf n_{ij}.
\end{equation}

The force is applied during both compression and rebound.  There is no
contact weighting, available-momentum cap, collision phase, or stored
internal momentum.  Starting overlaps are treated as ordinary physical
overlaps in this baseline model.

\VerAForceFunctionSection{GetContactStiffness}

\texttt{GetContactStiffness(SourceID, TargetID, contact\_type)} selects
particle-pair or wall stiffness for the current contact.

\VerAForceFunctionSection{GetPairStiffness}

\texttt{GetPairStiffness(SourceID, TargetID)} returns the nonnegative mean of
the two particle-owned stiffness values:

\begin{equation}
q_{ij}=\max\left(0,\frac{q_i+q_j}{2}\right).
\end{equation}

\VerAForceFunctionSection{CalcVelocity}

\texttt{CalcVelocity(SourceID, totalForce)} performs the source velocity
write.  It converts the completed source-local force sum to acceleration
inside the function:

\begin{equation}
\mathbf v_{n+1}
=
\mathbf v_n+\frac{\mathbf F_n}{m}\Delta t.
\end{equation}

\VerAForceFunctionSection{GetParticleMass}

\texttt{GetParticleMass(ParticleID)} returns particle mass from
\texttt{parms.x} with an epsilon lower bound.

\VerAForceFunctionSection{CalcPosition}

\texttt{CalcPosition(SourceID)} performs the semi-implicit Euler position
write using the newly calculated velocity:

\begin{equation}
\mathbf x_{n+1}
=
\mathbf x_n+\mathbf v_{n+1}\Delta t.
\end{equation}

This velocity-first update is the symplectic Euler method.  Symplectic Euler
does not conserve the system's exact total energy at every timestep.  Instead,
it typically produces a bounded energy excursion in which the calculated
total energy oscillates above and below the exact value.  Unlike standard
Euler integration, this error generally does not accumulate into continuous
energy growth or decay.  Decreasing the timestep reduces the excursion's
magnitude, making symplectic Euler well suited for long-running conservative
dynamics simulations.

\subsection{Potential Energy and Diagnostics}

The force is conservative with pair potential:

\begin{equation}
U_{ij}(d)=q_{ij}\int_d^{r_i+r_j}A_{ij}(s)\,ds.
\end{equation}

Potential energy, relative normal velocity, contact slots, end-position
snapshots, and pair-phase measurements are diagnostic only.  They do not
affect the source-local force sum or particle integration.  System energy is:

\begin{equation}
E_{\mathrm{total}}=KE+\sum_{i<j}U_{ij}+\sum_{i,w}U_{iw}.
\end{equation}

\subsection{Current Numerical Results}

The central pair rule independently produces equal-and-opposite particle-pair
forces and preserves vector momentum.  Semi-implicit Euler introduces a small
bounded numerical energy excursion.  Current captures include:

\begin{verbatim}
ThreeParticle, 500 steps:
energy range = approximately [-9.18e-5, +9.98e-5]
final energy difference = approximately -1.69e-7

ParticlesInBox, 10000 steps:
energy range = approximately [-1.85e-5, +1.90e-5]
final energy difference = approximately -7.34e-7
\end{verbatim}

The long-running box result currently rises and falls within a bounded band
rather than showing monotonic energy drift.

